\BellaFrontTitle{Avant-propos}
\origpage{XI}
Les enfants nous poussent souvent à entreprendre. C'est pour mon fils aîné, et en pensant au cadet, que j'ai commencé à rédiger ces notes, et, de page en page, je me suis mis à l'ouvrage pour obtenir un exposé d'une construction possible des Mathématiques.

Partant de la théorie des ensembles, l'introduction des cardinaux permet de construire $\mathbb N$, ensemble des entiers naturels. L'étude des structures de groupe conduit à symétriser $\mathbb N$ pour obtenir l'ensemble $\mathbb Z$ des entiers relatifs. La structure d'anneau et de corps donnera alors le corps des rationnels.

C'est dans le deuxième tome de cet ouvrage que j'aborderai la construction du corps des réels, de nature topologique. Mais dans celui-ci, après l'étude des espaces vectoriels et des polynômes, on construira le corps des complexes, $\mathbb C$, clôture algébrique de $\mathbb R$.

Tout au long de cette construction se dégage l'importance de la structure quotient attachée aux relations d'équivalence. L'étude du calcul matriciel et de l'algèbre extérieure me permet ensuite de traiter la réduction des endomorphismes. Enfin, les formes quadratiques et les formes hermitiennes, abordées ici d'un point de vue algébrique, se retrouveront en analyse dans l'étude des espaces préhilbertiens réels et dans celle des séries de Fourier.

Je me suis efforcé de rédiger ces notes pour les étudiants, en pensant à leurs réactions, qui figurent les miennes, face à une question ardue. Aussi pourra-t-on me reprocher un certain manque d'application, de ci de là, dans la formulation abstraite, l'utilisation des quantificateurs, mais j'ai voulu écrire un manuel vivant.

Je remercie mon épouse et mes enfants qui m'ont supporté pendant ces longues heures, mais aussi mes élèves du Lycée Saint-Louis qui, par leurs questions, m'ont conduit à préciser mes propres connaissances, au fil des années. Je remercie également Monsieur Cau, du service de reprographie du lycée, sans la gentillesse et la compétence duquel ces notes n'auraient \BellaCorrection{peut-être pas été mises en forme}{peut être pas été mis en formes}.

Enfin, je suis très reconnaissant envers Monsieur Deheuvels qui m'a encouragé à passer de notes rédigées pour des élèves à un traité plus élaboré.

\begin{flushright}Dun-les-Places, le 23 février 1993\end{flushright}
\clearpage
